\documentclass[11pt]{article}
\usepackage[margin=1in]{geometry}
\usepackage{amsmath,amssymb}
\usepackage[hidelinks]{hyperref}

\newcommand{\rad}{\operatorname{rad}}
\newcommand{\linf}{\ell_\infty}

\begin{document}

\section*{Human Review: \texorpdfstring{$c_0\subset\ell_\infty$}{c0 in ell-infinity} closed-bounded radii}

\noindent Original automatic proof: \texttt{solution\_packet.pdf}

\noindent Checked date: 10 June 2026

\noindent Status: Manually checked.

\noindent Verdict: Correct, but not a full solution to the stated characterisation problem.

\section*{Human Comments}

The argument is correct. The example \(c_0\subset\ell_\infty\) shows that
equality of radii for finite subsets does not extend automatically to all
closed bounded subsets.

However, this should not be read as a real solution to the broader problem as
stated, since the problem asks for characterisations of when equality holds.
The packet instead gives a natural obstruction: even when the relevant
finite-dimensional or finite-subset conditions hold, they may fail to control
the behaviour of closed bounded subsets.

In the example, every finite subset \(F\subset c_0\) has
\[
  \rad_{c_0}(F)=\rad_{\ell_\infty}(F),
\]
because a centre in \(\ell_\infty\) can be truncated to an element of \(c_0\)
without losing the radius bound on finitely many vectors. But for
\[
  A=\{e_n:n\in\mathbb N\}\subset c_0,
\]
one has
\[
  \rad_{\ell_\infty}(A)=\frac12,\qquad \rad_{c_0}(A)=1.
\]
The missing centre is the constant sequence
\[
  (1/2,1/2,\ldots),
\]
which belongs to \(\ell_\infty\) but not to \(c_0\).

\section*{Short Version}

This is a correct and useful counterexample to a possible extension from
finite subsets to closed bounded subsets. It shows that the finite behaviour
is not enough to determine the closed-bounded behaviour.

It does not provide a full characterisation of those pairs or subspaces for
which the radius equality holds for all closed bounded subsets.

\section*{Possible Follow-Up}

It may be worth checking with the authors whether they were already aware of
this example, and whether they consider it interesting enough to include as a
basic obstruction to the closed-bounded version of the problem.

\section*{Review Outcome}

Archive this packet as correct, with the caveat that its contribution is best
described as a natural counterexample or obstruction, not as a full solution
of the characterisation problem.

\end{document}
